\documentclass{article}
\usepackage{amsmath}
\usepackage{amssymb} % Add this package for \mathbb{}

\begin{document}

\title{Types of Numbers}
\author{}
\date{}
\maketitle

\section{Introduction}
In mathematics, there are several different types of numbers, each with its own properties and applications. These number systems form the foundation for more advanced mathematical concepts. This report outlines the main sets of numbers, including imaginary and complex numbers, and how they relate to each other.

\section{Natural Numbers (\( \mathbb{N} \))}
The set of natural numbers includes the counting numbers, starting from 1, and sometimes 0. The set can be defined as:
\[
\mathbb{N} = \{0, 1, 2, 3, 4, 5, \dots\}
\]
- \textbf{Properties}: Natural numbers are used for counting and ordering.
- \textbf{Note}: Some definitions exclude 0 from \( \mathbb{N} \), starting from 1.

\section{Whole Numbers}
The set of whole numbers includes the natural numbers, but always starts from 0. It is defined as:
\[
\text{Whole Numbers} = \{0, 1, 2, 3, 4, \dots\}
\]
- \textbf{Properties}: Whole numbers are natural numbers with the inclusion of 0.

\section{Integers (\( \mathbb{Z} \))}
The set of integers includes all positive and negative whole numbers, as well as zero. It is represented as:
\[
\mathbb{Z} = \{\dots, -3, -2, -1, 0, 1, 2, 3, \dots\}
\]
- \textbf{Properties}: Integers are used in scenarios where both positive and negative values are needed, such as in financial calculations (debt) or temperature.

\section{Rational Numbers (\( \mathbb{Q} \))}
The set of rational numbers includes all numbers that can be expressed as a fraction of two integers:
\[
\mathbb{Q} = \left\{\frac{a}{b} \mid a \in \mathbb{Z}, b \in \mathbb{Z}, b \neq 0 \right\}
\]
- \textbf{Properties}: Rational numbers include integers, terminating decimals, and repeating decimals.
- \textbf{Example}: \( \frac{1}{2}, -\frac{3}{4}, 0.75 \)

\section{Irrational Numbers}
Irrational numbers are real numbers that cannot be expressed as a ratio of two integers. Their decimal expansions are non-terminating and non-repeating.
\[
\text{Irrational Numbers} = \{ x \mid x \notin \mathbb{Q} \}
\]
- \textbf{Properties}: Examples include numbers like \( \pi \) and \( \sqrt{2} \), which cannot be exactly represented as a fraction.
- \textbf{Examples}: \( \pi = 3.14159\ldots \), \( \sqrt{2} = 1.41421\ldots \)

\section{Real Numbers (\( \mathbb{R} \))}
The set of real numbers includes all rational and irrational numbers. Real numbers represent any point along the number line.
\[
\mathbb{R} = \{ x \mid x \text{ is rational or irrational} \}
\]
- \textbf{Properties}: Real numbers include integers, rational numbers, and irrational numbers.
- \textbf{Examples}: \( -1, 0, 1.5, \sqrt{3}, \pi \)

\section{Imaginary Numbers}
Imaginary numbers are numbers that involve the square root of negative numbers. The fundamental unit of imaginary numbers is \( i \), which is defined as:
\[
i = \sqrt{-1}
\]
- \textbf{Properties}: Imaginary numbers are multiples of \( i \), and they exist outside of the real number line.
- \textbf{Examples}: \( 2i, -5i \)

\section{Complex Numbers (\( \mathbb{C} \))}
Complex numbers consist of two parts: a real part and an imaginary part. A complex number is written in the form:
\[
z = a + bi
\]
Where:
- \( a \) is the real part,
- \( b \) is the coefficient of the imaginary part, and
- \( i = \sqrt{-1} \) is the imaginary unit.

- \textbf{Properties}: Complex numbers extend the real number system to include solutions to equations like \( x^2 = -1 \), which have no real solutions.
- \textbf{Examples}: \( 3 + 4i, -1 + i, 0 + 5i \)

\section{Other Mainstream Number Sets}

\subsection{Prime Numbers}
Prime numbers are natural numbers greater than 1 that have no divisors other than 1 and themselves. The set of prime numbers is infinite.
\[
\text{Prime Numbers} = \{2, 3, 5, 7, 11, 13, 17, \dots\}
\]

\subsection{Transcendental Numbers}
Transcendental numbers are real or complex numbers that are not solutions to any non-zero polynomial equation with rational coefficients.
- \textbf{Examples}: \( \pi \), \( e \)

\subsection{Algebraic Numbers}
Algebraic numbers are any numbers that are solutions to polynomial equations with rational coefficients.
- \textbf{Example}: \( \sqrt{2} \) is algebraic because it is the solution to \( x^2 - 2 = 0 \).

\section{Conclusion}
Understanding the different types of numbers—natural, integer, rational, real, imaginary, and complex—is fundamental to exploring various mathematical concepts. Each type of number plays a specific role in calculations and theoretical mathematics, and they often work together in more advanced mathematical fields like algebra, calculus, and number theory.

\end{document}
